\clearpage
\renewcommand{\Article}[1]{\par\medskip\Needspace{4\baselineskip}\pdfbookmark[1]{Art. #1}{nexus-second-prior-art-#1}\centerline{\textbf{#1.}}\nobreak\smallskip}
\pdfbookmark[0]{Klassenanzahl und Determinante. II}{nexus-second-prior}
\begin{center}
[II.]\par\medskip
\Large\bfseries ÜBER DEN ZUSAMMENHANG ZWISCHEN DER ANZAHL DER KLASSEN, IN DIE BINÄRE QUADRATISCHE FORMEN ZERFALLEN, UND IHRER DETERMINANTE.\par\medskip
\normalsize ERSTE ABHANDLUNG\par\medskip
DER KÖNIGLICHEN GESELLSCHAFT VORGELEGT 1837\ .\ .\ .
\end{center}
\BookRule{.18\textwidth}

\Article{1}
Sechsunddreißig Jahre sind vergangen, seit die Prinzipien des bemerkenswerten Zusammenhangs entdeckt wurden, der in dieser Abhandlung behandelt werden soll, wie schon am Schluss der \emph{Untersuchungen über höhere Arithmetik} angekündigt wurde. Andere Beschäftigungen hatten die Aufmerksamkeit jedoch lange von dieser Untersuchung abgelenkt, bis es in jüngerer Zeit gelang, zu ihr zurückzukehren und sie durch neue Bemühungen zu erweitern. Da aber dieser neue Teil der höheren Arithmetik den Umfang einer einzigen Abhandlung übersteigt, soll diese erste den Formen mit negativen Determinanten gewidmet sein; die Formen mit positiven Determinanten, die eine ganz eigene Behandlung erfordern, bleiben einer anderen Abhandlung vorbehalten.

\SourceLeaf{287}
\Article{2}
Für unser Vorhaben wird ein Satz benötigt, der an sich zwar arithmetisch ist, dessen Wesen sich jedoch bequemer und deutlicher durch Überlegungen in geometrischer Form vor Augen führen lässt.

Ist in einer unbegrenzten Ebene eine durch eine beliebige Linie begrenzte Figur gegeben, so lässt sich ihr Flächeninhalt näherungsweise bestimmen, indem man die Ebene in Quadrate teilt und sowohl diejenigen zählt, die ganz innerhalb der Figur liegen, als auch diejenigen, die von ihrem Rand geschnitten werden. Offenbar ergibt sich ein kleinerer oder größerer Flächeninhalt als der wahre, je nachdem, ob die letzteren Quadrate weggelassen oder zu den ersteren hinzugezählt werden. Will man aber die letzteren, auf dem Rand liegenden Quadrate nach irgendeiner Regel teils ausschließen, teils mitzählen, so kann der Fehler positiv oder negativ sein; er muss jedoch kleiner sein als der gesamte Flächeninhalt aller Randquadrate. Je kleiner man die Quadrate wählt, desto genauer wird der Flächeninhalt auf diese Weise bestimmt. Man kann eine solche Annäherung unbegrenzt fortsetzen oder die Quadrate so klein wählen, dass der Fehler kleiner als jede gegebene Größe wird. Obwohl dies schon an sich einleuchtend erscheinen mag, werden wir es dennoch nicht verschmähen, es durch einen strengen Beweis abzusichern.

Zwei Quadrate können entweder einen Eckpunkt, zwei oder keinen gemeinsam haben; im ersten und zweiten Fall sollen sie benachbart, im dritten disjunkt heißen. Offenbar können höchstens vier Quadrate sämtlich einander benachbart sein; daher müssen sich unter je fünf verschiedenen Quadraten mindestens zwei disjunkte finden. Da nun der Abstand zwischen zwei Punkten in disjunkten Quadraten nicht kleiner sein kann als die Seitenlänge der Quadrate, die wir mit $a$ bezeichnen wollen, ist klar: Geht ein Punkt von einer beliebigen Stelle eines Quadrats aus, durchquert dann ein zweites, drittes und viertes und gelangt schließlich zu einem fünften, so ist die Länge seines Weges sicher nicht kleiner als $a$. Aus demselben Grund kann, wenn die Linie immer weitere Quadrate durchläuft, der Teil zwischen dem fünften und neunten sowie zwischen dem neunten und dreizehnten usw. nicht kleiner als $a$ sein. Wir schließen daher leicht, dass eine beliebige in sich zurückkehrende Linie, die insgesamt $n$ verschiedene Quadrate berührt hat, sicher nicht kürzer als $\frac{(n-4)a}4$ sein kann. Umgekehrt kann also eine geschlossene Linie der Länge $=l$ sicher nicht mehr als $4+\frac{4l}a$ verschiedene Quadrate berührt haben. Da deren Flächeninhalt $=4aa+4al$ bei unbegrenzt abnehmendem $a$ kleiner als jede gegebene Größe werden kann, gilt dies umso mehr für den Fehler der oben besprochenen Flächenbestimmung.

\SourceLeaf{288}
\Article{3}
Die Regel für die Zulassung oder den Ausschluss der auf dem Rand der Figur liegenden Quadrate ließe sich auf viele verschiedene Arten festlegen. Am einfachsten scheint es jedoch, nur die Lage des Mittelpunkts jedes Quadrats zu beachten: Quadrate, deren Mittelpunkte innerhalb der Figur liegen, werden zugelassen, diejenigen, deren Mittelpunkte außerhalb liegen, ausgeschlossen. Schließlich bleibt es unserer Wahl überlassen, ob Mittelpunkte, die etwa auf dem Rand selbst liegen, den inneren oder den äußeren zugerechnet werden sollen. Anstelle der Mittelpunkte könnten wir auch beliebige andere Punkte nehmen, die in den einzelnen Quadraten jeweils dieselbe relative Lage haben.

Damit kommt es darauf an, sich in der Ebene Punkte in gleichen Abständen auf ebenfalls gleich weit voneinander entfernten Geraden so verteilt zu denken, dass sie Quadrate bilden. Nach dem Satz des vorigen Artikels können wir dann behaupten, dass die Anzahl der in der Figur enthaltenen Punkte, multipliziert mit dem Quadrat des Abstands zweier nächster Punkte, dem Flächeninhalt der Figur beliebig nahe kommt, sofern jener Abstand nur hinreichend klein gewählt wird; oder, nach der gewöhnlichen Ausdrucksweise, dass dieses Produkt den Flächeninhalt angibt, wenn der Abstand unendlich klein ist.

\Article{4}
Die durch die folgende Gleichung zwischen den rechtwinkligen Koordinaten $p,q$ ausgedrückte Kurve,
\[
app+2bpq+cqq=1
\]
ist ein Kegelschnitt, und zwar eine Ellipse, wenn $a,c$ sowie $ac-bb$ positive Größen sind. Der von dieser Ellipse eingeschlossene Flächeninhalt ergibt sich zu $\pdot=\frac\pi{\surd(ac-bb)}$. Der Wert von $app+2bpq+cqq$ ist außerhalb der Ellipse überall größer als $1$, innerhalb kleiner als $1$ und nirgends negativ.

Man denke sich ein .System von Punkten in der Ebene der Ellipse so verteilt, dass sie Quadrate bilden, deren Seiten $=\lambda$ den Koordinatenachsen parallel sind. Dabei ist es gleichgültig, ob der Koordinatenursprung, also der Mittelpunkt der Ellipse, mit einem dieser Punkte zusammenfällt oder nicht. Die Anzahl der Punkte innerhalb der Ellipse, einschließlich etwa auf dem Rand selbst liegender Punkte, sei $=m$. Dann ist nach dem Satz des vorigen Artikels $\frac\pi{\surd(ac-bb)}$ der Grenzwert von $m\lambda\lambda$, dem diese Größe bei unbegrenzt abnehmendem $\lambda$ beliebig nahe kommt.

Nehmen wir an, dass der Koordinatenursprung mit einem Punkt des Systems zusammenfällt, so sind, wenn wir $p=\lambda x,q=\lambda y$ setzen, offenbar für die einzelnen Punkte des Systems $x$ und $y$ ganze Zahlen; umgekehrt entspricht jede Kombination ganzzahliger Werte der Größen
\SourceLeaf{289}
$x,y$ einem Punkt des Systems. Die Zahl $m$ ist daher nichts anderes als die Anzahl aller Kombinationen ganzzahliger Werte von $x,y$, für die $F$ nicht größer als $M$ wird, wenn wir der Kürze wegen die Funktion oder quadratische Form $axx+2bxy+cyy$ mit $F$ und die Größe $\frac1{\lambda\lambda}$ mit $M$ bezeichnen. Die Determinante dieser Form ist $bb-ac$; dafür schreiben wir $-D$. Unser Satz ist nun folgendermaßen auszusprechen.

\textsc{Satz I.}\quad\emph{Die Anzahl $m$ aller Kombinationen ganzzahliger Werte der Unbestimmten $x,y$, für die der Wert einer Form mit negativer Determinante $-D$ die Grenze $M$ nicht überschreitet, wird $=\frac{\pi M}{\surd D}$, zwar näherungsweise, doch mit unbegrenzt zunehmender Annäherung, während $M$ ins Unendliche wächst. Es braucht kaum darauf hingewiesen zu werden, dass die unbegrenzte Annäherung hier, und ebenso im Folgenden, nicht so zu verstehen ist, als ob die Differenz zwischen $\frac{\pi M}{\surd D}$ und $m$ selbst unbegrenzt abnähme. Vielmehr nähert sich das Verhältnis dieser Größen unbegrenzt der Gleichheit; das heißt, $\frac{\pi M}{m\surd D}-1$ nimmt unbegrenzt ab.}

\Needspace{13\baselineskip}
\Article{5}
Um die Auszählung tatsächlich auszuführen, kann man für die einzelnen ganzzahligen Werte von $x$ zwischen den Grenzen $-\surd\frac{cM}D$ und $+\surd\frac{cM}D$ die beiden Werte von $y$ berechnen, die $F=M$ entsprechen; daraus ergibt sich unmittelbar die Anzahl der dazwischenliegenden ganzen Zahlen. Da diese Anzahl für entgegengesetzte Werte von $x$ dieselbe ist, ersparen wir uns nahezu die Hälfte der Arbeit. Man kann auch die Werte von $x$ zählen, die den einzelnen Werten von $y$ zwischen den Grenzen $-\surd\frac{aM}D$ und $+\surd\frac{aM}D$ entsprechen. Durch eine geeignete Verbindung beider Methoden lässt sich die Arbeit weiter verringern; wir führen dies hier jedoch nicht ausführlicher aus. Es möge genügen, einiges zum einfachsten Fall hinzuzufügen.

Die Form sei $F=xx+yy$, also die Kurve ein Kreis; ferner seien $r,r',r'',r'''\ldots r^{(r)}$ die nächstkleineren ganzen Zahlen unterhalb von
\[
\surd M,\quad\surd(M-1),\quad\surd(M-4),\quad\surd(M-9)\ldots\surd(M-rr)
\]
oder diese Größen selbst, falls sich unter ihnen ganze Zahlen befinden. Dann ist die gesuchte Anzahl
\[
\begin{aligned}
m&=2r+1+2(2r'+1)+2(2r''+1)+2(2r'''+1)+\text{usw.}+2(2r^{(r)}+1)\\
&=1+4r+4r'+4r''+4r'''+\text{usw.}+4r^{(r)}
\end{aligned}
\]

Dasselbe erreichen wir jedoch schneller, indem wir mit $q$ die nächstkleinere ganze Zahl
\SourceLeaf{290}
unterhalb von $\surd\tfrac12 M$ bezeichnen (oder diese Größe selbst, wenn sie ganzzahlig ist) und die Formel
\[
m=4qq+1+4r+8(r^{(q+1)}+r^{(q+2)}+r^{(q+3)}+\text{usw.}+r^{(r)})
\]
verwenden. Auf diese Weise wurden die folgenden Anzahlen ermittelt.
\[
\begin{array}{r|r||r|r||r|r}
M&m&M&m&M&m\\\hline
100&317&1000&3149&10000&31417\\
200&633&2000&6293&20000&62845\\
300&949&3000&9425&30000&94237\\
400&1257&4000&12581&40000&125629\\
500&1581&5000&15705&50000&157093\\
600&1885&6000&18853&60000&188453\\
700&2209&7000&21993&70000&219901\\
800&2521&8000&25137&80000&251305\\
900&2821&9000&28269&90000&282697\\
1000&3149&10000&31417&100000&314197
\end{array}
\]

\Needspace{14\baselineskip}
\Article{6}
Dem Satz des Artikels 4 geben wir auf folgende Weise größere Allgemeinheit.

\textsc{Satz II.}\quad\emph{Sollen nicht alle Kombinationen ganzzahliger Werte von $x,y$, für die $F$ den Wert $M$ nicht überschreitet, gesammelt werden, sondern nur diejenigen in bestimmten Abständen, nämlich diejenigen, bei denen $x$ einer gegebenen Zahl $G$ zu einem gegebenen Modul $g$ und $y$ einer gegebenen Zahl $H$ zu einem gegebenen Modul $h$ kongruent ist, so wird die Anzahl $m'$ dieser Kombinationen näherungsweise durch $\frac{\pi M}{gh\surd D}$ ausgedrückt, mit unbegrenzt zunehmender Annäherung, während $M$ ins Unendliche wächst.}

Setzen wir nämlich $x=gx'+G,y=hy'+H$, so ist klar, dass $m'$ die Anzahl aller Kombinationen ganzzahliger Werte von $x',y'$ ist, für die
\[
agg(x'+\tfrac Gg)^2+2bgh(x'+\tfrac Gg)(y'+\tfrac Hh)+chh(y'+\tfrac Hh)^2
\]
den Wert $M$ nicht überschreitet. Denken wir uns also ein Punktesystem in der Ebene wie in Artikel 4 verteilt, jedoch so, dass nicht der Koordinatenursprung, sondern der Punkt mit den Koordinaten $p=\frac{G\lambda}g,q=\frac{H\lambda}h$ mit einem Punkt des Systems zusammenfällt, so drückt $m'$ die Anzahl der Punkte innerhalb der
\SourceLeaf{291}
Ellipse aus, deren Gleichung
\[
aggpp+2bghpq+chhqq=1
\]
ist, wobei etwa auf dem Rand selbst liegende Punkte stets mitgezählt werden. Der Flächeninhalt dieser Ellipse $=\frac\pi{gh\surd(ac-bb)}=\frac\pi{gh\surd D}$ ist der Grenzwert, dem sich das Produkt $m'\lambda\lambda=\frac{m'}M$ unbegrenzt nähert, wenn $\lambda$ unbegrenzt abnimmt oder $M$ ins Unendliche wächst.

Offenbar umfasst unser Satz auch den Fall, dass nur eine der Unbestimmten $x,y$ in bestimmten Abständen fortschreiten soll, während der Wert der anderen keiner Bedingung unterliegt. Denn dies ist offenbar dasselbe, als wenn entweder $h$ oder $g$ $=1$ gesetzt wird.

\Article{7}
Die bisherigen Ausführungen sind von der Beschaffenheit der Koeffizienten der Form $axx+2bxy+cyy$ unabhängig; von nun an werden wir aber voraussetzen, dass diese Koeffizienten ganzzahlig sind. Damit gibt jede Kombination ganzzahliger Werte von $x,y$ der Form selbst einen ganzzahligen Wert oder entspricht einer Darstellung irgendeiner ganzen Zahl durch diese Form. Daher ist die Gesamtheit aller Kombinationen ganzzahliger Werte von $x,y$, durch die die Form $F=axx+2bxy+cyy$ einen Wert nicht größer als die Grenze $M$ erhält, dieselbe wie die Gesamtheit aller Darstellungen ganzer Zahlen, die die Grenze $M$ nicht überschreiten, also bis zu dieser Grenze einschließlich, falls sie selbst ganzzahlig ist. Bezeichnen wir daher \mbox{der Kürze wegen} die Anzahl der verschiedenen Darstellungen einer bestimmten ganzen Zahl $n$ durch die Form $F$ mit $F(n)$ oder, wenn keine Mehrdeutigkeit zu befürchten ist, einfach mit $Fn$, so ist die oben durch $m$ ausgedrückte Zahl $=F0+F1+F2+F3+\text{usw.}+FM$, und der erste Satz nimmt die folgende Gestalt an.

\textsc{Satz III.}\quad\emph{Die Summe $F0+F1+F2+\text{usw.}+FM$ wird näherungsweise durch $\frac{\pi M}{\surd D}$ ausgedrückt, mit unbegrenzt zunehmender Annäherung, während $M$ ins Unendliche wächst.}

\Article{8}
Es bietet sich an, dem dritten Satz, der die Darstellungen \emph{aller} Zahlen betrifft, einen anderen hinzuzufügen, der nur ungerade Zahlen betrifft. Offenbar können durch die Form $F$ keine ungeraden Zahlen dargestellt werden, wenn $a$ und $c$ beide gerade Zahlen sind; die Untersuchung wird daher auf die drei übrigen Fälle beschränkt.

I. Ist $a$ ungerade und $c$ gerade, so wird eine ungerade Zahl dargestellt, indem man $x$ einen ungeraden Wert gibt, während der Wert von $y$ beliebig bleibt. Satz II
\SourceLeaf{292}
zeigt daher, wenn man $g=2,G=1,h=1$ setzt, dass die Anzahl aller Kombinationen solcher Werte von $x,y$, die der Form einen ungeraden Wert nicht größer als die Grenze $M$ geben, mit unbegrenzter Annäherung durch $\frac{\pi M}{2\surd D}$ ausgedrückt wird, während $M$ ins Unendliche wächst.

II. Ist $a$ gerade und $c$ ungerade, so muss zur Darstellung einer ungeraden Zahl $y$ ungerade sein. Setzen wir daher $g=1,h=2,H=1$, so gelangen wir zum selben Schluss.

III. Sind sowohl $a$ als auch $c$ ungerade, so muss entweder ein ungerader Wert von $x$ mit einem geraden Wert von $y$ oder ein gerader Wert von $x$ mit einem ungeraden Wert von $y$ kombiniert werden, damit sich ein ungerader Wert der Formel ergibt. Für jede der beiden Arten wird die Anzahl aller Kombinationen, für die der Wert der Form die Grenze $M$ nicht überschreitet, mit unbegrenzter Annäherung durch $\frac{\pi M}{4\surd D}$ ausgedrückt. Daher wird die Anzahl aller Kombinationen, die der Form ungerade Werte nicht größer als die Grenze $M$ geben, auch hier mit unbegrenzter Annäherung durch $\frac{\pi M}{2\surd D}$ ausgedrückt.

Da nun die Gesamtheit aller solchen Kombinationen nichts anderes ist als die Gesamtheit aller Darstellungen aller Zahlen $1,3,5,7\ldots M$, wenn $M$ eine ungerade ganze Zahl ist, oder $1,3,5,7\dotsfour M-1$, wenn $M$ gerade ist, haben wir

\begingroup
\setlength{\abovedisplayskip}{3pt plus 1pt minus 1pt}
\setlength{\belowdisplayskip}{3pt plus 1pt minus 1pt}
\setlength{\abovedisplayshortskip}{3pt plus 1pt minus 1pt}
\setlength{\belowdisplayshortskip}{3pt plus 1pt minus 1pt}
\textsc{Satz IV.}\quad\emph{Die Summe}
\[
F1+F3+F5+F7\ldotp\ldotp+FM\quad\text{oder}\quad F1+F3+F5+F7\ldotp\ldotp+F(M-1)
\]
\emph{(je nachdem, ob $M$ ungerade oder gerade ist) wird mit unbegrenzter Annäherung durch $\frac{\pi M}{2\surd D}$ ausgedrückt, sofern $F$ eine Form ist, in der einer der Koeffizienten $a,c$ oder beide ungerade sind.}
\par\nobreak\begingroup\offinterlineskip\smallskip
\centerline{\rule{.30\textwidth}{.35pt}}\smallskip\endgroup
\endgroup
